% ===========================================================================
\section*{Open threads · where this goes next}
\addcontentsline{toc}{section}{Open threads · where this goes next}
\warmth{0}

\begin{strand}
These notes are an \keyword{arsenal}, not a museum. The cross-section open questions below are the
seeds for the next pass: a unifying perspective, the missing lower bounds, and the exit toward
empirical processes.
\end{strand}

\block{Open threads (cross-section)}
\begin{itemize}[leftmargin=1.4em,itemsep=2pt]
  \item \keyword{Four faces}: concentration $\Leftrightarrow$ isoperimetry $\Leftrightarrow$ log-Sobolev
        $\Leftrightarrow$ transportation (Marton). Drawing the explicit commuting diagram of these equivalences
        is the ``main theorem'' this notebook most wants to add.
        \loose{Once this diagram is in, §\ref{sec:entropy}–§\ref{sec:talagrand} unify into one engine.}
  \item \keyword{When does dimension enter}? Gaussian concentration is dimension-free, matrix concentration carries
        $\log d$, chaining carries metric entropy. A ``sources of dimension dependence'' table separates ``truly
        dimension-free'' from ``hidden in the constants''.
  \item \keyword{Tightness}: is every upper bound paired with a lower bound? Cramér (§\ref{sec:chernoff}) and
        Sudakov (§\ref{sec:chaining}) are; the lower bounds for McDiarmid and matrix Bernstein (worst $f$ / worst spectrum) are to be filled.
  \item \keyword{The exit}: after §\ref{sec:chaining}, open a dedicated \emph{empirical-process} notebook (VC / Rademacher / uniform laws / Donsker).
\end{itemize}

\bigskip
\begin{strand}
Closing line: the whole subject is seven variations on one move, \emph{exponentiate the deviation, then crush it
with structure}. What the structure is decides which engine you reach for. Recognizing the structure beats memorizing the formula.
\end{strand}
